\documentclass[11pt]{article}

\usepackage[margin=1in]{geometry}
\usepackage[T1]{fontenc}
\usepackage{lmodern}
\usepackage{microtype}
\usepackage{amsmath,amssymb,amsthm,mathtools}
\usepackage{booktabs}
\usepackage{array}
\usepackage{enumitem}
\usepackage[dvipsnames]{xcolor}
\usepackage[colorlinks=true,linkcolor=MidnightBlue,citecolor=BrickRed,urlcolor=MidnightBlue]{hyperref}
\usepackage[nameinlink,capitalise,noabbrev]{cleveref}
\usepackage{fancyhdr}
\usepackage{mdframed}
\usepackage{tikz}
\usetikzlibrary{arrows.meta,positioning,calc}

\setlength{\parindent}{0pt}
\setlength{\parskip}{0.65em}
\setlist{itemsep=0.25em,topsep=0.35em}
\pagestyle{fancy}
\fancyhf{}
\lhead{Mass Production and Local Sharing}
\rhead{Working research note}
\cfoot{\thepage}
\setlength{\headheight}{14pt}

\newtheorem{theorem}{Theorem}[section]
\newtheorem{lemma}[theorem]{Lemma}
\newtheorem{proposition}[theorem]{Proposition}
\newtheorem{corollary}[theorem]{Corollary}
\newtheorem{conjecture}[theorem]{Conjecture}
\theoremstyle{definition}
\newtheorem{definition}[theorem]{Definition}
\newtheorem{remark}[theorem]{Remark}
\newtheorem{example}[theorem]{Example}

\newcommand{\bits}{\{0,1\}}
\newcommand{\F}{\mathbb{F}}
\newcommand{\C}{\mathrm{C}}
\newcommand{\NCC}{\mathrm{NCC}}
\newcommand{\DT}{\mathrm{D}}
\newcommand{\Cert}{\mathrm{C}_1}
\newcommand{\poly}{\operatorname{poly}}
\newcommand{\supp}{\operatorname{supp}}
\newcommand{\Mux}{\operatorname{Mux}}
\newcommand{\RM}{\operatorname{RM}}
\newcommand{\wt}{\operatorname{wt}}
\newcommand{\eps}{\varepsilon}
\newcommand{\tC}{\widetilde O}
\newcommand{\xor}{\mathbin{\oplus}}
\newcommand{\bigxor}{\bigoplus}
\newcommand{\given}{\,|\,}

\title{\textbf{Mass Production and Local Sharing in Boolean Circuit Complexity}\\[0.4em]
\large An exponential-range construction and a sparse-share refinement of a meta-complexity reduction}
\author{Samuel Schlesinger\\\small Working research note developed through an exploratory collaboration with GPT-5.6 Pro}
\date{August 2026}

\begin{document}
\maketitle

\begin{mdframed}[backgroundcolor=Yellow!10,linecolor=BrickRed,linewidth=1pt]
\textbf{Status and caution.}
This is a working research note, not a peer-reviewed paper.  The central mass-production theorem below is accompanied by a self-contained proof that has been internally stress-tested, but its historical novelty remains unresolved.  In particular, the complete quantitative statements in the older simultaneous-realization literature of Paul, Uhlig, and Albrecht have not all been available for comparison.  The sparse-share modification of Hirahara's reduction is likewise presented as a candidate refinement that requires independent checking by an expert familiar with the original proof.  The document deliberately distinguishes proved implications of the displayed arguments from novelty claims.
\end{mdframed}

\begin{abstract}
A theorem of Uhlig says that an arbitrary Boolean function on $n$ bits can be evaluated on $2^{o(n/\log n)}$ independent inputs using essentially the worst-case circuit size required for one evaluation.  We develop a local view of this phenomenon in which sharing is measured by the number of calls to an opaque subcomputation.  Deterministic local sharing becomes decision-tree query complexity, while nondeterministic local sharing becomes certificate complexity.  This viewpoint exposes Uhlig's two-copy construction as a disjoint-recovery code and motivates a higher-dimensional construction.

The main candidate theorem states that for every fixed $\gamma<1$, every Boolean function $f:\bits^n\to\bits$, and every $t\le 2^{\gamma n}$, the direct product $f^{\times t}$ has circuit size $O_\gamma(2^n/n)$.  The proof uses a constant-rate low-degree evaluation code, affine-line recovery, deterministic disjoint scheduling inside groups, controlled congestion between groups, and a recursive bootstrap
\[
   \beta\longmapsto \frac{1}{2-\beta}.
\]
Starting from the disjoint range $\beta=1/2$, iteration reaches every constant exponent below one.  A consequence is that any circuit lower-bound measure dominated by circuit size and additive under independent copies is subexponential on every Boolean function.

We then audit an application of Uhlig's theorem in Hirahara's NP-hardness reduction for partial MCSP.  The stronger mass-production range is not the active parameter bottleneck there.  The audit instead identifies a separate quadratic loss: padding a length-$\Delta$ share for every one of up to $\Delta$ parties.  A leafwise formula secret-sharing scheme has only $\Delta$ share bits in total.  With occurrence-indexed Nisan--Wigderson masks and a regularizing tag, the sampler has input length $O(\log n+\Delta)$ rather than $O(\log n+\Delta^2)$ and is exactly uniform over a support of density $2^{-\Delta}$.  Re-running the average-case density argument shows that Hirahara's AveMCSP conversion remains valid with $\beta=2^{-\Delta-3}$, after harmless padding.  Tracking the same CMMSA exponent, one may take $\Delta=\Theta(\log n)$ and improve the explicit partial-MCSP gap exponent from $a/2$ to $a$, without changing the qualitative $(\log N)^{\Omega(1)}$ statement.
\end{abstract}

\tableofcontents
\newpage

\section{Overview}

For a Boolean function $f:\bits^n\to\bits$ and an integer $t\ge 1$, write
\[
 f^{\times t}(x^{(1)},\ldots,x^{(t)})
   =\bigl(f(x^{(1)}),\ldots,f(x^{(t)})\bigr).
\]
The naive upper bound is
\[
 \C(f^{\times t})\le t\C(f).
\]
For unrestricted Boolean circuits, this can be very far from tight.  Uhlig proved that if
$t=2^{o(n/\log n)}$, then every $f$ admits a circuit for $f^{\times t}$ of size
$(1+o(1))2^n/n$ and depth $(1+o(1))n$; see \cite{Uhlig1974,Uhlig1992,Wegener1987,Kretschmer2023}.  Thus a mildly subexponential number of evaluations can be performed for essentially the Shannon--Lupanov cost of one worst-case function.

The motivating observation of this note is that the sharing in such a circuit need not be a globally reusable preprocessing phase.  A subcomputation may be useful only to a particular coalition of inputs, and which coalition uses it may depend on the actual input values.  The right primitive is therefore \emph{conditional demand}: a costly module can be shared whenever its possible uses have low simultaneous congestion.

The work has four components.

\begin{enumerate}[label=\textbf{\arabic*.}]
\item We define parametric local reductions in which an arbitrary module $G$ is treated as opaque.  The minimum number of $G$-occurrences is exactly deterministic oracle-query complexity.  With existential witnesses, it is exactly nondeterministic query complexity, equivalently $1$-certificate complexity.

\item We identify Uhlig's two-copy decomposition with a primitive two-request batch-code gadget.  This suggests replacing a single telescoping path by a constant-rate evaluation code with many affine-line recovery sets.

\item We give a candidate proof that, for every fixed $\gamma<1$,
\[
   t\le 2^{\gamma n}
   \quad\Longrightarrow\quad
   \C(f^{\times t})=O_\gamma(2^n/n)
\]
for every $f$.  The construction does not require a globally disjoint schedule.  It enforces disjointness only within request groups and recursively mass-produces the resulting bounded number of repeated resource evaluations.

\item We re-run the parameters in Hirahara's reduction \cite{Hirahara2022}.  The improved mass-production range does not directly strengthen the published result, because the old Uhlig range already covers the required number of calls.  A different optimization appears available: use one share bit per formula leaf occurrence and index pseudorandom masks by variable-occurrence pairs.  We present a regular sampler, a blockwise hybrid argument, and the corresponding AveMCSP density calculation.
\end{enumerate}

\section{Circuit conventions and direct products}

We work over a fixed finite complete basis of fan-in at most two, such as
$\{\wedge,\vee,\neg\}$.  Circuit size is the number of non-input gates.  Fan-out and output wires are free.  Multi-output circuits are allowed.  Replacing one finite complete basis by another changes all size bounds by only a constant factor.

Let
\[
 L(n)=\max_{f:\bits^n\to\bits}\C(f).
\]
The Shannon counting lower bound and Lupanov's upper bound give
\[
 L(n)=\Theta(2^n/n),
\]
and the sharper asymptotic constant is known for standard bases.  Our main construction uses only the order $O(2^n/n)$; unlike Uhlig's sharp result, it does not presently preserve the leading constant.

For later induction, define
\[
 L_r(n)=\max_{f:\bits^n\to\bits}\C(f^{\times r}).
\]
The target statement is that $L_r(n)=O_\gamma(2^n/n)$ whenever
$r\le 2^{\gamma n}$ and $\gamma<1$ is fixed.

\section{Local sharing as oracle-query complexity}

\subsection{Parametric circuit contexts}

Let $A$ be a finite address space, and let
\[
 G:A\to\bits
\]
be an uninterpreted module.  A \emph{parametric circuit context} $H[G]$ is a Boolean circuit containing ordinary Boolean gates and occurrences of the gate $G$.  The ordinary gates may compute addresses for later occurrences of $G$ from the external input and from earlier answers.  We count only occurrences of $G$.

For external input $x$, the context defines a functional
\[
 \Phi_x:\bits^A\to\bits^m,
 \qquad
 \Phi_x(G)=H[G](x).
\]
Let $D(\Phi_x)$ denote deterministic decision-tree complexity when a query asks for one truth-table bit $G(a)$.

\begin{theorem}[Local sharing equals query complexity]\label{thm:oracle-query}
The minimum number of occurrences of the opaque module $G$ in a parametric circuit computing $\Phi_x(G)$ for every $x$ and every interpretation of $G$ is
\[
 \max_x D(\Phi_x).
\]
\end{theorem}

\begin{proof}
Topologically order the $G$-gates in a circuit containing $q$ occurrences.  After fixing $x$, the address supplied to the first occurrence is determined.  The address supplied to the second can depend on the first answer, and so on.  The circuit therefore induces an adaptive decision tree of depth at most $q$ for $\Phi_x$.

Conversely, let $q=\max_xD(\Phi_x)$.  For every external input $x$, fix a depth-$q$ decision tree for $\Phi_x$, padding shorter paths.  Since the surrounding Boolean logic is free in this resource measure, it can compute the next address as an arbitrary function of $x$ and the previous answers, and can compute the final output from the transcript.  This compiles the family of trees into $q$ sequential occurrences of $G$.
\end{proof}

This exact equivalence is useful because it turns a vague notion of ``sharing an internal computation'' into a familiar combinatorial invariant.

\subsection{Demand width}

Suppose a context has guarded potential calls
\[
  p_i(x)\cdot G(a_i(x)),\qquad i\in[m],
\]
and must expose every guarded answer.  Define the set of distinct simultaneously demanded addresses
\[
  \mathcal D(x)=\{a_i(x):p_i(x)=1\}
\]
and the demand width
\[
  \Delta_G=\max_x|\mathcal D(x)|.
\]

\begin{corollary}[Black-box demand theorem]\label{cor:demand}
For an arbitrary module $G$, the exact minimum number of physical copies required to expose all guarded answers is $\Delta_G$.
\end{corollary}

\begin{proof}
The upper bound queries each distinct active address once and fans the answer out.  For the lower bound, fix an input attaining $\Delta_G$.  If fewer addresses are queried, choose two interpretations of $G$ agreeing on all queried addresses and differing on one unqueried demanded address.
\end{proof}

The key lesson is that anti-sharing lower bounds must establish unavoidable \emph{simultaneous demand}.  Functional independence of subcomputations alone is insufficient: an entirely arbitrary $G$ can still be shared if its demands are mutually exclusive.

A particularly simple universal reduction is
\[
 \Mux_c(G(a),G(b))=G(\Mux_c(a,b)),
\]
where the mux on the right is applied coordinatewise.  Two possible calls to an arbitrary module collapse to one whenever only the selected answer is needed.

\subsection{Nondeterminism and certificates}

Now assume the output is Boolean and allow an existential witness:
\[
 \Phi_x(G)=1
 \quad\Longleftrightarrow\quad
 \exists w\;V(x,w;G)=1.
\]
Let $C_1(\Phi_x)$ denote the maximum, over positive truth tables of $G$, of the minimum number of queried truth-table positions needed to certify that $\Phi_x(G)=1$.

\begin{theorem}[Nondeterministic local sharing equals certificate complexity]\label{thm:oracle-cert}
The minimum number of occurrences of $G$ in an existential verifier for $\Phi$ is
\[
  \max_x C_1(\Phi_x).
\]
\end{theorem}

\begin{proof}
An accepting computation path of a verifier reveals a set of oracle answers whose values force acceptance; this is a $1$-certificate.  Conversely, guess a minimum certificate and its claimed values, query the listed addresses, and verify both the answers and the certificate relation.  As in \cref{thm:oracle-query}, arbitrary dependence on $x$ is absorbed into the free surrounding logic.
\end{proof}

For example,
\[
 \Phi(G)=G(a_1)\vee\cdots\vee G(a_m)
\]
has deterministic query complexity $m$ in the worst case but $1$-certificate complexity one.  In this local model, nondeterminism is literally the ability to guess which possible resource demand is the successful one.

\section{Uhlig's two-copy construction as disjoint recovery}

Fix a prefix length $p$, put $N=2^p$, and define restrictions
\[
 f_i(z)=f(i,z),\qquad i\in\{0,\ldots,N-1\}.
\]
Uhlig's two-copy construction introduces
\[
 g_0=f_0,\qquad
 g_i=f_{i-1}\xor f_i\quad(1\le i<N),\qquad
 g_N=f_{N-1}.
\]
Then every $f_i$ has two representations:
\[
 f_i=\bigxor_{\ell=0}^{i}g_\ell
     =\bigxor_{\ell=i+1}^{N}g_\ell.
\]
For requests $i\le j$, recover $f_i$ from the prefix set
$\{0,\ldots,i\}$ and $f_j$ from the suffix set
$\{j+1,\ldots,N\}$.  These sets are disjoint, so every expensive function $g_\ell$ is evaluated on at most one suffix.

This is precisely a linear two-request disjoint-recovery code.  Up to an invertible linear change of coordinates, it is the familiar length-$N+1$ single-parity-check code, the basic primitive two-multiset batch code in the terminology of Ishai, Kushilevitz, Ostrovsky, and Sahai \cite{IKOS2004}.

The coding viewpoint suggests a general template.  Encode all prefix restrictions pointwise in the suffix.  A requested restriction is recovered from a set of code coordinates.  If the recovery sets assigned to simultaneous requests are disjoint, each coordinate function is evaluated only once.  If they have congestion $C$, each coordinate function needs at most $C$ evaluations.  The next sections construct a code with many line recoveries and then recursively absorb this congestion.

\section{A constant-rate low-degree resource code}

Fix an integer $\ell\ge2$ and a field $\F_q$ of characteristic two, where $q=2^b$.  Let
\[
 s_0=\left\lfloor\frac{q-1}{\ell}\right\rfloor
\]
and choose a set $A\subseteq\F_q$ of size $s_0$.  Consider the vector space $\mathcal P$ of polynomials
\[
 P(X_1,\ldots,X_\ell)\in\F_q[X_1,\ldots,X_\ell]
\]
whose degree in each variable is less than $s_0$.  Values on the grid $A^\ell$ determine such a polynomial uniquely.  Moreover,
\[
 \deg P\le \ell(s_0-1)\le q-2.
\]

The evaluation map
\[
 \operatorname{Enc}:\F_q^{A^\ell}\to\F_q^{\F_q^\ell},
 \qquad
 (P(u))_{u\in A^\ell}\mapsto(P(y))_{y\in\F_q^\ell},
\]
is therefore a systematic linear code of length
\[
 N=q^\ell
\]
and dimension
\[
 K=s_0^\ell=\Theta_\ell(q^\ell).
\]
Its rate is a positive constant depending only on $\ell$.

\subsection{Packing the restrictions}

Split an $n$-bit input as
\[
 x=(a,z),\qquad |a|=p,\quad |z|=d=n-p.
\]
For a fixed suffix $z$, pack the $2^p$ bits
\[
 (f(a,z))_{a\in\bits^p}
\]
into $K$ field symbols, using $b$ bits per symbol and zero-padding unused positions.  Choose $q$ minimally, among powers of two, so that
\[
 Kb\ge 2^p.
\]
For fixed $\ell$, minimality gives
\begin{equation}\label{eq:packing}
 q^\ell b=\Theta_\ell(2^p).
\end{equation}
Let $P_z$ be the interpolating polynomial.  For every code coordinate
$y\in\F_q^\ell$, define the field-valued resource function
\[
 G_y(z)=P_z(y).
\]
Each of the $b$ output bits of $G_y$ is an arbitrary Boolean function on $d$ variables.

\subsection{Affine-line recovery}

\begin{lemma}[Line identity]\label{lem:line-identity}
Let $P:\F_q^\ell\to\F_q$ be represented by a polynomial of total degree at most $q-2$.  For every $u\in\F_q^\ell$ and every nonzero direction $v\in\F_q^\ell$,
\[
 P(u)=-\sum_{\lambda\in\F_q^\ast}P(u+\lambda v).
\]
In characteristic two the minus sign may be omitted.
\end{lemma}

\begin{proof}
The restriction $Q(\lambda)=P(u+\lambda v)$ is a univariate polynomial of degree at most $q-2$.  For every monomial $\lambda^j$ with $0\le j\le q-2$, the sum over all $\lambda\in\F_q$ is zero.  Hence
$\sum_{\lambda\in\F_q}Q(\lambda)=0$, which gives the identity after separating the term $\lambda=0$.
\end{proof}

Thus every affine line through an information point $u\in A^\ell$ supplies a recovery set consisting of its $q-1$ non-target points.  The number of projective directions through a point is
\begin{equation}\label{eq:directions}
 D_q=\frac{q^\ell-1}{q-1}=q^{\ell-1+o(1)}.
\end{equation}
Repeated requests for the same information point pose no special problem: two distinct lines through that point intersect only at the omitted target.

\section{Deterministic scheduling with controlled congestion}

\subsection{Disjoint scheduling inside one group}

Consider a group of $g$ requests, represented by target information points
$u_1,\ldots,u_g\in A^\ell$.  Process them in order.  Suppose recovery sets have already been selected for requests $1,\ldots,j-1$, and let $U$ be their union.  Because the sets are disjoint,
\[
 |U|=(j-1)(q-1)<gq.
\]
For the current target $u_j$, each point $y\in U\setminus\{u_j\}$ forbids at most one projective direction, namely the direction of $y-u_j$.  The point $u_j$ itself forbids none, since the target is omitted from its own recovery set.  We obtain the following.

\begin{lemma}[Greedy disjoint line scheduler]\label{lem:greedy-scheduler}
If
\[
  gq<D_q,
\]
then every multiset of $g$ target points admits pairwise disjoint line recovery sets.  The lines can be selected by a deterministic Boolean circuit of size
\[
 \widetilde O_\ell(g^2q),
\]
where the tilde hides factors polynomial in $\log q$ and $\log g$.
\end{lemma}

\begin{proof}
The counting argument above guarantees an unused direction at every step.  For constructivity, canonically rank projective directions by normalizing the first nonzero coordinate to one.  At step $j$, compute and sort the $O(jq)$ forbidden ranks, deduplicate them, and choose the first missing rank.  Standard sorting and priority circuits use $\widetilde O_\ell(jq)$ gates.  Enumerating the $q-1$ points of the selected line has the same order of cost.  Summing over $j\le g$ gives $\widetilde O_\ell(g^2q)$.
\end{proof}

\subsection{Grouping and congestion}

Partition $t$ requests into $C$ fixed groups, each of size at most
\[
 g=\left\lceil\frac{t}{C}\right\rceil.
\]
Apply \cref{lem:greedy-scheduler} independently within every group.  A code coordinate is used at most once per group, hence at most $C$ times overall.  The total scheduling size is
\begin{equation}\label{eq:scheduler-cost}
 C\cdot\widetilde O_\ell(g^2q)
   =\widetilde O_\ell\!\left(\frac{t^2q}{C}\right).
\end{equation}
The construction is valid when
\begin{equation}\label{eq:scheduler-condition}
 \frac{tq}{C}<D_q.
\end{equation}

The suffix records can be routed to slots indexed by
$(\text{group},\text{code coordinate})$.  Inside a group there is at most one record per coordinate.  Sorting/scatter networks and inverse routing use
\begin{equation}\label{eq:routing-cost}
 \widetilde O_\ell(tq+Cq^\ell)
\end{equation}
gates, including payload widths polynomial in $n$.  There is no hidden $tq^\ell$ crossbar: one sorts the $tq$ actual incidence records together with one dummy record per one of the $Cq^\ell$ slots.

\section{The exponential-range mass-production theorem}

\subsection{A master recurrence}

Let
\[
 L_C(d)=\max_{h:\bits^d\to\bits}\C(h^{\times C}).
\]
There are $q^\ell$ resource coordinates and $b$ Boolean output bits per coordinate.  Since every coordinate has congestion at most $C$, the resource layer costs
\[
 q^\ell b\,L_C(d).
\]
Combining the resource layer, scheduling, routing, and line decoding gives the following bound.

\begin{proposition}[Master bound]\label{prop:master}
Let $n=p+d$, choose $q$ by \cref{eq:packing}, and suppose \cref{eq:scheduler-condition} holds.  Then every $f:\bits^n\to\bits$ satisfies
\begin{equation}\label{eq:master}
 \C(f^{\times t})
 \le
 q^\ell b\,L_C(d)
 +\widetilde O_\ell\!\left(
      \frac{t^2q}{C}+tq+Cq^\ell
   \right).
\end{equation}
\end{proposition}

\begin{proof}
For each target prefix, the scheduler chooses a recovery line.  Every line point $y$ requests the resource $G_y$ on the suffix belonging to that output instance.  The per-coordinate congestion is at most $C$, so each output bit of $G_y$ is evaluated by a circuit for its $C$-fold direct product.  The resulting field values are routed back and XORed along each line using \cref{lem:line-identity}; the requested bit is selected from the recovered information symbol.  Equations \eqref{eq:scheduler-cost} and \eqref{eq:routing-cost} account for all non-resource gates.
\end{proof}

\subsection{The inductive property}

For $0<\beta\le1$, let $\mathcal P(\beta)$ denote the assertion:

\begin{quote}
For every fixed $\rho<\beta$, there is a constant $A_\rho$ such that, for every sufficiently large $m$, every $h:\bits^m\to\bits$, and every $r\le2^{\rho m}$,
\[
 \C(h^{\times r})\le A_\rho\frac{2^m}{m}.
\]
\end{quote}

\subsection{The base range}

\begin{lemma}[Base case]\label{lem:base}
$\mathcal P(1/2)$ holds.
\end{lemma}

\begin{proof}
Fix $\gamma<1/2$ and let $t\le2^{\gamma n}$.  Put
\[
 p=(1-\gamma)n+O(1),\qquad d=\gamma n+O(1),\qquad C=1.
\]
By \cref{eq:packing},
\[
 \log_2 q=\frac{(1-\gamma)n}{\ell}+o(n),
 \qquad
 \log_2 q^\ell=(1-\gamma)n+o(n).
\]
The scheduling condition has exponent inequality
\[
 \gamma+\frac{1-\gamma}{\ell}
 <(1-\gamma)\left(1-\frac1\ell\right),
\]
which holds for sufficiently large fixed $\ell$ because $\gamma<1/2$.

The scheduler term has exponent
\[
 2\gamma+\frac{1-\gamma}{\ell}<1
\]
for the same choice of $\ell$.  The terms $tq$ and $q^\ell$ also have exponents strictly below one.  Finally, by \cref{eq:packing} and the Shannon--Lupanov upper bound,
\[
 q^\ell b\,L(d)
 =O_\ell\!\left(2^p\frac{2^d}{d}\right)
 =O_\gamma(2^n/n).
\]
Every remaining term is $2^{(1-\Omega_\gamma(1))n}\poly(n)$ and is therefore negligible relative to $2^n/n$.
\end{proof}

\subsection{The bootstrap}

\begin{lemma}[Congestion bootstrap]\label{lem:bootstrap}
If $\mathcal P(\beta)$ holds, then
\[
 \mathcal P\!\left(\frac{1}{2-\beta}\right)
\]
holds.
\end{lemma}

\begin{proof}
Fix
\[
 \gamma<\frac1{2-\beta}.
\]
The inequality is equivalent to
\[
 2\gamma-1<\beta\gamma.
\]
The range $\gamma\le1/2$ is already covered by \cref{lem:base}, so assume $\gamma>1/2$.  Choose a constant $c$ satisfying
\begin{equation}\label{eq:c-window}
 2\gamma-1<c<\beta\gamma.
\end{equation}
Put
\[
 p=(1-\gamma)n+O(1),\qquad d=\gamma n+O(1),
 \qquad C=2^{cn+O(1)}.
\]
Choose $\rho$ with $c/\gamma<\rho<\beta$.  Then
\[
 C\le2^{\rho d},
\]
so $\mathcal P(\beta)$ gives
\[
 L_C(d)=O_\rho(2^d/d).
\]
Consequently, the resource term in \cref{eq:master} is again $O_\gamma(2^n/n)$.

Let $\theta=1-\gamma$.  The group size has exponent $\gamma-c$.  The scheduling condition becomes
\[
 \gamma-c+\frac{\theta}{\ell}
 <\theta\left(1-\frac1\ell\right).
\]
That is,
\[
 c>2\gamma-1+\frac{2\theta}{\ell}.
\]
This holds for sufficiently large fixed $\ell$ by the left inequality in \eqref{eq:c-window}.

The three overhead exponents in \cref{eq:master} are
\begin{align*}
 \frac1n\log_2\!\left(\frac{t^2q}{C}\right)
   &=2\gamma-c+\frac\theta\ell+o(1)<1,\\
 \frac1n\log_2(tq)
   &=\gamma+\frac\theta\ell+o(1)<1,\\
 \frac1n\log_2(Cq^\ell)
   &=c+\theta+o(1)<1.
\end{align*}
The first follows from the stronger scheduling inequality, the second from $\ell>1$, and the third from $c<\beta\gamma\le\gamma$.  Hence every overhead term is exponentially below $2^n/n$.
\end{proof}

\begin{theorem}[Exponential-range Boolean mass production]\label{thm:main-mass}
For every fixed $\gamma<1$, there is a constant $A_\gamma$ such that for every sufficiently large $n$, every Boolean function $f:\bits^n\to\bits$, and every
$t\le2^{\gamma n}$,
\[
 \C(f^{\times t})\le A_\gamma\frac{2^n}{n}.
\]
\end{theorem}

\begin{proof}
Starting from $\beta_1=1/2$, define
\[
 \beta_{k+1}=\frac1{2-\beta_k}.
\]
Induction gives
\[
 \beta_k=\frac{k}{k+1}.
\]
By \cref{lem:base,lem:bootstrap}, $\mathcal P(\beta_k)$ holds for every $k$.  Given fixed $\gamma<1$, choose $k$ with $\gamma<k/(k+1)$.
\end{proof}

\begin{remark}[Quantitative comparison with Uhlig]
The previously quoted uniform range is
\[
 t=2^{o(n/\log n)}
\]
with a sharp $(1+o(1))2^n/n$ leading term.  \Cref{thm:main-mass} permits
\[
 t=2^{(1-\eps)n}
\]
for every fixed $\eps>0$, but currently only with a constant $A_\eps$ that can deteriorate rapidly as $\eps\to0$.  The improvement is therefore in the exponent of the number of copies, not in the leading constant.
\end{remark}

\begin{remark}[Why the bootstrap matters]
A globally disjoint schedule is unnecessary.  Disjointness inside $C$ groups creates congestion $C$ across groups.  The induction hypothesis evaluates those $C$ repeated resource calls for the same asymptotic cost as one.  In slogan form,
\[
 \text{global sharing}
 =\text{local disjoint sharing}+\text{recursive reuse across layers}.
\]
\end{remark}

\section{Consequences and barriers}

\subsection{A barrier to copy-additive lower-bound measures}

\begin{corollary}[Additive-measure barrier]\label{cor:additive-barrier}
Suppose $\mu$ is a nonnegative complexity measure on Boolean vector functions satisfying
\[
 \mu(F)\le\C(F)
\]
and
\[
 \mu(f^{\times r})\ge r\mu(f)
\]
for every $f$ and $r$.  Then, for every fixed $\eps>0$ and every $f:\bits^n\to\bits$,
\[
 \mu(f)=O_\eps(2^{\eps n}/n).
\]
In particular, $\mu(f)=2^{o(n)}$ uniformly over all Boolean functions.
\end{corollary}

\begin{proof}
Apply \cref{thm:main-mass} with $r=\lfloor2^{(1-\eps)n}\rfloor$:
\[
 r\mu(f)\le\mu(f^{\times r})\le\C(f^{\times r})
   \le A_{1-\eps}\frac{2^n}{n}.
\]
Divide by $r$.
\end{proof}

Thus a lower-bound potential that assigns independent hardness to independent copies cannot prove exponential unrestricted-circuit lower bounds.  Any successful measure must account for the possibility that the copies reorganize their computation through cross-instance sharing.

\subsection{An amortization curve}

It is natural to refine the scalar $\C(f)$ to the curve
\[
 A_f(\gamma,n)=\C\!\left(f^{\times\lfloor2^{\gamma n}\rfloor}\right).
\]
The main theorem gives the universal upper envelope
\[
 A_f(\gamma,n)=O_\gamma(2^n/n)
 \qquad(\gamma<1).
\]
Restricted models may have very different curves.  Formulas, bounded fan-out circuits, monotone circuits, branching programs, and bounded-depth circuits become natural test cases for a theory of shareability.  The local demand viewpoint suggests that the relevant obstruction is not simply functional hardness but unavoidable congestion under the routing and selection operations available in the model.

\section{Auditing the Hirahara application}

Hirahara's reduction from CMMSA to partial meta-complexity uses Uhlig's theorem in Lemma~8.3 of \cite{Hirahara2022}.  We first determine whether \cref{thm:main-mass} directly improves that reduction.

Let $n_0$ be the number of CMMSA variables, let $\Delta$ be the formula degree, and let $\eps_0$ be the CMMSA soundness.  Hirahara sets
\[
 \delta=\frac1{\log n_0},
 \qquad
 \eps=\frac{\eps_0}{2\Delta^2}.
\]
Each amplified value is locally computable with
$O(1/(\eps\delta))$ nonadaptive calls to an underlying random function $f_k$, and the completeness circuit needs at most $\Delta$ amplified values per selected variable.  The number of calls mass-produced for one $f_k$ is therefore
\begin{equation}\label{eq:hirahara-old-q}
 Q_{\mathrm{old}}
 =O\!\left(\frac{\Delta}{\eps\delta}\right)
 =O\!\left(\frac{\Delta^3\log n_0}{\eps_0}\right).
\end{equation}

The upstream CMMSA theorem supplies constants $a,\beta>0$ such that one may take
\[
 \eps_0=\Delta^{-a},
 \qquad
 \Delta\le2^{(\log n_0)^{1-\beta}}.
\]
Even at this largest permitted degree,
\[
 \log Q_{\mathrm{old}}
 =O((\log n_0)^{1-\beta})
 =o\!\left(\frac{\log n_0}{\log\log n_0}\right).
\]
The truth-table input length of $f_k$ is
$u=\Theta(\log\lambda)=\Theta(\log n_0)$ for the polynomial-size instances in the reduction.  Hence
\[
 Q_{\mathrm{old}}=2^{o(u/\log u)},
\]
which already lies in Uhlig's classical range.  We conclude:

\begin{proposition}[The new range is not the active Hirahara bottleneck]\label{prop:hirahara-not-active}
Directly replacing Uhlig's theorem by \cref{thm:main-mass} does not enlarge the degree range available in Hirahara's published parameter optimization.  The original number of calls is already admissible under the classical $2^{o(u/\log u)}$ theorem throughout the upstream CMMSA range.
\end{proposition}

The active loss is elsewhere.  Hirahara gives each of at most $\Delta$ parties a share padded to length $\Delta$, producing a target input of length
\[
 O(\log n_0+\Delta^2).
\]
Since the reduction explicitly enumerates the target truth table, polynomial running time forces $\Delta^2=O(\log n_0)$ and hence $\Delta=O(\sqrt{\log n_0})$.

\section{Leafwise secret sharing for monotone formulas}

Let $\varphi$ be a monotone formula with $L$ literal occurrences.  Binarize all gates without changing the leaves.  Let $A(\varphi)$ and $O(\varphi)$ be the numbers of binary AND and OR gates.  Then
\begin{equation}\label{eq:tree-count}
 A(\varphi)+O(\varphi)=L-1.
\end{equation}

Define the following linear secret-sharing scheme over $\F_2$.

\begin{itemize}
\item Assign the root the secret $b$.
\item At an OR gate, assign both children the parent's value.
\item At an AND gate with parent value $c$, choose a fresh uniform bit $r$, assign $r$ to the left child and $c\xor r$ to the right child.
\item Give each leaf's arriving bit to the party labelling that occurrence.  A repeated variable receives one bit for each occurrence.
\end{itemize}

\begin{lemma}[Leafwise formula sharing]\label{lem:leaf-sharing}
The construction is a perfect secret-sharing scheme for the access structure represented by $\varphi$.  It uses exactly one share bit per literal occurrence, so the total share length is $L$.  Reconstruction from every authorized set is linear over $\F_2$.
\end{lemma}

\begin{proof}
Correctness follows by induction.  At an OR gate, an authorized set satisfies at least one child and recovers the common parent value.  At an AND gate, it satisfies both children and XORs the recovered values $r$ and $c\xor r$ to obtain $c$.

For privacy, induct on the formula.  At an OR gate, an unauthorized set is unauthorized in both subformulas; the two subviews use independent internal randomness and are each independent of the common secret.  At an AND gate, the set is unauthorized in at least one child.  The random split $r$ one-time-pads the secret seen by either child alone, while the unauthorized child's view is independent of its assigned value by induction.  Therefore the joint view is independent of the parent secret.

Every reconstruction operation is XOR or projection, hence linear.
\end{proof}

\begin{lemma}[Injectivity of the full share map]\label{lem:share-injective}
Let $R\in\bits^{A(\varphi)}$ be the vector of AND-gate coins and let
\[
 S_\varphi:\bits^{1+A(\varphi)}\to\bits^L
\]
map $(b,R)$ to the full vector of leaf shares.  Then $S_\varphi$ is injective.
\end{lemma}

\begin{proof}
Given all leaf shares, recursively recover the value assigned to every node.  At an OR gate the parent equals either child value.  At an AND gate, both child values determine the parent by XOR and also reveal the split coin, taken to be the left-child value.  Thus all of $(b,R)$ is uniquely determined.
\end{proof}

This injectivity is not needed for secret sharing itself, but it is crucial for making the modified Hirahara sampler uniform over its support.

\section{A sparse and regular version of Hirahara's sampler}

This section gives the promised line-by-line modification.  We retain Hirahara's hardness amplification, Kolmogorov-complexity argument, and Nisan--Wigderson extraction lemma, changing only the share layout and the indexing of repeated pseudorandom coordinates.

\subsection{Dyadic formula indices}

A finite fair-coin circuit samples a dyadic distribution.  To avoid an inessential exact-sampling issue, choose
$h_\Phi=\lceil\log_2\nu\rceil$ and a balanced map
\[
 \pi:\bits^{h_\Phi}\to[\nu]
\]
whose fibers differ in size by at most one.  Include the raw index $J\in\bits^{h_\Phi}$ in the target input and use formula $\varphi_{\pi(J)}$.  Every formula then has probability within a factor two of $1/\nu$.  Replacing the upstream soundness parameter by a constant-factor smaller value absorbs this distortion.  To simplify notation below, we write $\varphi_J$ for the selected formula and treat $J$ as the formula index.

\subsection{Occurrence-indexed NW masks}

Construct one NW design
\[
 \{S_{k,r}:k\in[n_0],\ r\in[\Delta]\}
\]
of $n_0\Delta$ sets.  For a formula $\varphi_J$, enumerate its leaves
$q=1,\ldots,L_J$.  Let $k_J(q)$ be the variable labelling leaf $q$, and let $r_J(q)$ be the rank of that occurrence among the leaves labelled by $k_J(q)$.  The $q$th leaf share is masked by
\[
 \widehat f_{k_J(q)}(z_{S_{k_J(q),r_J(q)}}).
\]
This global indexing fixes a subtle bookkeeping point: the same repeated function $\widehat f_k$ always occupies the block of coordinates $(k,1),\ldots,(k,\Delta)$, independently of the formula in which it appears.

\subsection{Regularizing the sampler}

Let
\[
 L_J=L(\varphi_J),\qquad A_J=A(\varphi_J),\qquad O_J=O(\varphi_J).
\]
The sampler chooses independently and uniformly:
\[
 b\in\bits,
 \quad R\in\bits^{A_J},
 \quad u\in\bits^{\Delta-L_J},
 \quad v\in\bits^{O_J},
 \quad z\in\bits^d.
\]
Let $\sigma=S_{\varphi_J}(b,R)\in\bits^{L_J}$ be the leaf-share vector.  Define two $\Delta$-bit blocks $\xi$ and $\tau$ by
\begin{align*}
 \xi_q&=
 \widehat f_{k_J(q)}(z_{S_{k_J(q),r_J(q)}})\xor\sigma_q,
 &&1\le q\le L_J,\\
 (\xi_{L_J+1},\ldots,\xi_\Delta)&=u,\\
 (\tau_1,\ldots,\tau_{O_J})&=v,\\
 (\tau_{O_J+1},\ldots,\tau_\Delta)&=0^{\Delta-O_J}.
\end{align*}
The target input and label are
\begin{equation}\label{eq:sparse-sample}
 x=(J,z,\xi,\tau),
 \qquad\text{label }b.
\end{equation}
The tag $\tau$ is independent of the secret and is ignored by the completeness algorithm.  Its purpose is purely regularity.

By \cref{eq:tree-count}, the number of random bits other than $J$ and $z$ is
\begin{align*}
 1+A_J+(\Delta-L_J)+O_J
 &=1+(L_J-1-O_J)+(\Delta-L_J)+O_J\\
 &=\Delta.
\end{align*}

\begin{proposition}[Exact support regularity]\label{prop:regularity}
For every fixed $J$ and $z$, the map
\[
 (b,R,u,v)\longmapsto(\xi,\tau)
\]
is injective and has image size $2^\Delta$.  Consequently the sampler in \eqref{eq:sparse-sample} is uniform over its support $D$, and if
\[
 m_0=h_\Phi+d+2\Delta
\]
is the target input length, then
\begin{equation}\label{eq:density}
 |D|=2^{h_\Phi+d+\Delta}=2^{m_0-\Delta},
 \qquad
 \frac{|D|}{2^{m_0}}=2^{-\Delta}.
\end{equation}
Moreover, no input receives conflicting labels.
\end{proposition}

\begin{proof}
The tag reveals $v$, the padded part of $\xi$ reveals $u$, and the first $L_J$ coordinates of $\xi$ reveal the complete leaf-share vector after XORing the fixed NW mask determined by $J,z$ and the reduction's hardwired functions.  By \cref{lem:share-injective}, the leaf shares uniquely determine $(b,R)$.  Thus the map is injective.  Its domain has $2^\Delta$ elements.  Since $J,z$ and all internal coins are uniform, every support point has the same probability.  Injectivity also makes $b$ a function of $x$ on the support.
\end{proof}

The input length is
\[
 |x|=O(\log n_0+\Delta),
\]
because the NW seed length is $O(\log n_0)$ under Hirahara's choice of $\lambda$.

\section{Sparse-share refinement of Hirahara's Lemma 8.3}

We state the refinement with constant factors in the soundness parameter suppressed; the dyadic formula indexing above costs at most a factor two and can be absorbed by starting from a correspondingly smaller CMMSA soundness.

\begin{theorem}[Candidate sparse-share refinement]\label{thm:sparse-hirahara}
In the setting of Hirahara's Lemma~8.3, replace the party-padded sampler by \eqref{eq:sparse-sample} and set
\[
 \delta=\frac1{\log n_0},
 \qquad
 \eps=\frac{\eps_0}{2\Delta}.
\]
Then the completeness and soundness conclusions of the lemma continue to hold, up to absolute constant rescaling of $\eps_0$, while
\[
 |x|=O(\log n_0+\Delta).
\]
The number of nonadaptive calls to each underlying function $f_k$ required by the completeness circuit is
\begin{equation}\label{eq:sparse-query-count}
 Q_{\mathrm{sparse}}
 =O\!\left(\frac{\Delta}{\eps\delta}\right)
 =O\!\left(\frac{\Delta^2\log n_0}{\eps_0}\right).
\end{equation}
The sampler is uniform over a support of density exactly $2^{-\Delta}$ before optional padding.
\end{theorem}

\subsection{Completeness}

Suppose $T\subseteq[n_0]$ has weight at most $\theta$ and satisfies every formula.  Hardwire $T$ and the strings $f_k$ for $k\in T$.  For every selected $k$, compute all $\Delta$ values
\[
 \widehat f_k(z_{S_{k,1}}),\ldots,
 \widehat f_k(z_{S_{k,\Delta}}).
\]
Whenever a leaf $q$ is labelled by $k\in T$, unmask $\sigma_q$ from $\xi_q$.  Because $T$ is authorized for $\varphi_J$, the recursive reconstruction of \cref{lem:leaf-sharing} outputs $b$.  The tag $\tau$ and padded coordinates are ignored.

Each amplified value uses $O(1/(\eps\delta))$ calls to $f_k$, giving \eqref{eq:sparse-query-count}.  At the final choice $\Delta=\Theta(\log n_0)$ and $\eps_0=\Delta^{-a}$, this is polylogarithmic and hence well inside the original Uhlig range.  Hirahara's remaining size, running-time, and depth estimates are unchanged or improve: formula lookup, leaf routing, and linear reconstruction have polynomial size and logarithmic depth, while the direct-product computation of the $\widehat f_k$ values remains the bottleneck.

\subsection{A blockwise form of the NW extraction step}

Hirahara applies his Lemma~6.6 to $\Delta$ repeated copies of every $\widehat f_k$.  For precision, it is useful to state the blockwise consequence explicitly.

\begin{lemma}[Repeated-function block hybrid]\label{lem:block-hybrid}
Let $g_1,\ldots,g_n:\bits^\ell\to\bits$, and for each $k$ let
$S_{k,1},\ldots,S_{k,R}$ be NW design sets with pairwise intersections at most $\rho$.  Let a distinguisher $D(\alpha,Z,(X_{k,r})_{k,r})$ depend, for every fixed advice $\alpha$, on at most $m'$ individual coordinates $X_{k,r}$.  Then for every $\eta>0$ there is a set $B\subseteq[n]$ such that:

\begin{enumerate}[label=(\roman*)]
\item replacing every entire block $(X_{k,1},\ldots,X_{k,R})$ with uniform bits for $k\notin B$ changes the acceptance probability by at most $\eta$, for every advice $\alpha$;
\item for every $k\in B$,
\[
 K^D_{1/2-\eta/(2m')}(g_k)
 \le nR\,2^\rho+|\alpha|+O(\log(nR|\alpha|d\ell'/\eta)),
\]
where $\ell'$ is a standard upper bound on the time-bounded description complexity of the $g_k$'s.
\end{enumerate}
\end{lemma}

\begin{proof}[Proof sketch]
Repeat the hybrid proof of Hirahara's Lemma~6.6, but define $B$ at the level of base functions: $k\in B$ exactly when $g_k$ has a description meeting the displayed uniform threshold.  If the blockwise hybrid gap exceeded $\eta$, telescope only over the at most $m'$ individual coordinates on which the fixed-advice distinguisher depends.  Some transition would predict one coordinate $g_k(Z_{S_{k,r}})$ with advantage at least $\eta/(2m')$.  The standard NW reconstruction encodes $g_k$ using the other oracle strings, the advice, and at most
$nR2^\rho$ overlap information, contradicting $k\notin B$.  Because the threshold is uniform over all $r$ in a block, the hybrid may keep or replace the whole block at once.
\end{proof}

This formulation avoids an otherwise ambiguous projection from a subset of the $n\Delta$ repeated coordinates to a subset of the $n$ underlying variables.

\subsection{Soundness}

Fix a program $M$ of size at most Hirahara's soundness threshold.  Define a distinguisher whose advice contains
\[
 \alpha=(b,J,\sigma,u,v,\text{the leaf labels and formula shape}),
\]
padded to a common length.  Given the NW seed $Z$ and candidate blocks
$Y_k\in\bits^\Delta$, it constructs
\[
 x=(J,Z,\xi,\tau)
\]
by using coordinate $Y_{k_J(q),r_J(q)}$ at leaf $q$, and accepts exactly when $M(x)=b$.
For fixed advice, the distinguisher depends on only
\[
 L_J\le\Delta
\]
NW coordinates.  Apply \cref{lem:block-hybrid} with $m'=\Delta$ and hybrid error $\eta=\Theta(\eps_0)$.  We obtain a set $B\subseteq[n_0]$ such that every $k\in B$ satisfies
\[
 K^D_{1/2-\eps_0/(2\Delta)}(\widehat f_k)=o(\lambda).
\]
Hirahara's hardness-amplification lemma, now with
$\eps=\eps_0/(2\Delta)$, implies
\[
 K^D(f_k)=o(w(k)\lambda)
 \qquad(k\in B).
\]
The same summation and incompressibility argument as in the original proof yields
\[
 w(B)<\theta.
\]
Hence, in a No instance, the set $B$ satisfies only an $O(\eps_0)$ fraction of the formulas under the dyadic formula distribution.

In the hybrid experiment, every block $Y_k$ for $k\notin B$ is uniform.  Therefore each masked leaf coordinate belonging to a party outside $B$ is itself uniform and its share can be removed from the view.  If $\varphi_J(B)=0$, the remaining leaf shares are precisely a subvector of the shares held by the unauthorized party set $B\cap V_J$.  By \cref{lem:leaf-sharing}, their joint distribution is independent of $b$.  The padding $u$ and tag $v$ are independently random and do not affect this conclusion.  Thus the hybrid success probability is exactly $1/2$ on unsatisfied formulas.  Adding the fraction of satisfied formulas and the NW hybrid error gives the same $1/2+O(\eps_0)$ soundness as Hirahara's argument.

This completes the proof audit of \cref{thm:sparse-hirahara}, conditional only on the cited hardness-amplification and NW extraction lemmas.

\section{Re-running the AveMCSP density argument}

Hirahara's Theorem~8.7 converts the partial-function reduction to AveMCSP by randomly completing the function outside its support.  The proof requires three properties: a lower and upper bound on the support density, enough ambient points for a Chernoff plus union bound, and uniformity of the sampler over its support.  The regular sampler was designed precisely to make these properties transparent.

\subsection{Density before padding}

By \cref{prop:regularity}, with input length $m_0$,
\[
 |D|=2^{m_0-\Delta}.
\]
Set
\begin{equation}\label{eq:new-beta}
 \beta=2^{-\Delta-3}.
\end{equation}
Then
\[
 \beta 2^{m_0+3}=|D|,
 \qquad
 |D|\le2^{m_0-1}
\]
for $\Delta\ge1$.  This is the exact analogue of Hirahara's original choice
$\beta=2^{-\Delta^2-3}$, with $\Delta^2$ replaced by $\Delta$.

\subsection{Harmless uniform padding}

The union bound also requires
\[
 s'\log s'=o(\beta^2 2^m),
\]
where $s'$ is the largest No-case circuit size over which the random completion is union-bounded.  If this is not already true, append $h$ independent uniform bits $r$ to the target input and define the label to ignore them.  The new support is
\[
 D'=D\times\bits^h,
\]
so the density remains $2^{-\Delta}$ and the distribution remains uniform.

Soundness is preserved: if a circuit $C(x,r)$ has high success over the padded distribution, averaging over $r$ yields a fixed $r_0$ for which $C(x,r_0)$ has at least the same success over the unpadded distribution.  Completeness circuits simply ignore $r$.

For example, take
\begin{equation}\label{eq:padding-choice}
 h=3\left\lceil\log_2(s'+2)\right\rceil+2\Delta+10.
\end{equation}
Then, writing $m=m_0+h$,
\[
 \beta^2 2^m
 \ge16(s'+2)^3,
\]
so $s'\log s'=o(\beta^2 2^m)$.  In Hirahara's parameter range $s'=\poly(n_0)$, and for $\Delta=O(\log n_0)$ the padding satisfies
$h=O(\log n_0)$; the full truth table remains polynomially large.

\begin{proposition}[AveMCSP conversion survives sparsification]\label{prop:ave}
With the sparse regular sampler, Hirahara's Chernoff and union-bound proof of AveMCSP NP-hardness applies unchanged after replacing
\[
 \beta=2^{-\Delta^2-3}
 \quad\text{by}\quad
 \beta=2^{-\Delta-3}
\]
and, if necessary, adding the padding in \eqref{eq:padding-choice}.
\end{proposition}

\begin{proof}
The three hypotheses used in the original proof are now explicit:
\begin{enumerate}[label=(\arabic*)]
\item $\beta2^{m+3}=|D'|\le2^{m-1}$;
\item $s'\log s'=o(\beta^22^m)$ by \eqref{eq:padding-choice};
\item the sample distribution is uniform over $D'$ by \cref{prop:regularity} and uniform padding.
\end{enumerate}
Randomly fill the complement of $D'$.  In the Yes case, a fixed completeness circuit agrees on $D'$ and obtains the required global advantage by Chernoff concentration on the complement.  In the No case, each circuit of size at most $s'$ has only the permitted advantage on $D'$, and its agreement on the random complement concentrates around one half.  The failure probability for one circuit is $\exp(-\Omega(\beta^22^m))$; union-bounding over $2^{O(s'\log s')}$ circuits succeeds by item~(2).
\end{proof}

\subsection{The resulting tracked gap}

Use the same CMMSA theorem as Hirahara, with
\[
 \eps_0=\Delta^{-a},
 \qquad
 g=\Delta^a.
\]
The sparse target length and optional padding are $O(\log n_0+\Delta)$.  We may therefore take
\[
 \Delta=\Theta(\log n_0)
\]
while keeping the target truth table polynomial in $n_0$.  If $N_{\mathrm{tt}}=2^m$ is that truth-table length, then
\[
 \log N_{\mathrm{tt}}=\Theta(\log n_0),
\]
and the tracked gap becomes
\[
 g=\Delta^a=\Theta((\log N_{\mathrm{tt}})^a).
\]
Hirahara's original substitution $\Delta=(\log n_0)^{1/2}$ gives the tracked exponent $a/2$.  Thus the sparse-share modification doubles the exponent inherited from the same CMMSA theorem:
\[
 (\log N_{\mathrm{tt}})^{a/2}
 \quad\longrightarrow\quad
 (\log N_{\mathrm{tt}})^a.
\]
The published qualitative headline remains $(\log N)^{\Omega(1)}$, because it suppresses the constant exponent.  The AveMCSP conclusion remains NP-hardness; this calculation does not by itself assert a new AveMCSP approximation factor.

\section{What is established by the argument, and what remains}

\subsection{Mathematical claims supported by the displayed proofs}

Subject to ordinary proof checking, the note supports the following claims.

\begin{enumerate}
\item Opaque-module multiplicity is exactly deterministic query complexity, and its existential version is certificate complexity.
\item The grouped affine-line construction obeys the master recurrence \eqref{eq:master}.
\item The recurrence gives the bootstrap
$\beta\mapsto1/(2-\beta)$ and hence \cref{thm:main-mass} for every fixed exponent below one.
\item Copy-additive lower-bound measures dominated by circuit size are necessarily subexponential.
\item The stronger mass-production range is not the active constraint in Hirahara's published optimization.
\item Leafwise formula sharing, occurrence-indexed NW masks, and the regularity tag reduce the target input length to $O(\log n_0+\Delta)$ while making the support density exactly $2^{-\Delta}$.
\item The blockwise hybrid and privacy argument preserve Hirahara's soundness up to constant rescaling, and the AveMCSP density proof survives after harmless padding.
\end{enumerate}

\subsection{Open verification and novelty tasks}

The highest-priority remaining tasks are:

\begin{enumerate}
\item Compare \cref{thm:main-mass} with the full theorems of Paul \cite{Paul1976}, Uhlig \cite{Uhlig1974,Uhlig1992}, and Albrecht \cite{Albrecht1989}.  The accessible Albrecht preview says that an ``extremely large'' number of inputs can be handled with moderate hardware growth, but does not expose the quantitative theorem.

\item Audit the scheduler at the gate-basis level, including a fully explicit projective-direction ranking and the sorting/scatter implementation.  No asymptotic obstruction is visible, but this is the most engineering-like part of the proof.

\item Check the sparse-share modification against the exact machine model and all constants in Hirahara's Lemmas~6.6, 8.1, and 8.3.  In particular, a specialist should verify the blockwise version of the repeated-function hybrid and the description-length bookkeeping.

\item Determine whether the regularizing-tag device, or an equivalent regular sampler, is already implicit in the secret-sharing or meta-complexity literature.

\item Optimize the constant $A_\gamma$ and ask whether the leading $(1+o(1))$ Shannon--Lupanov constant can be preserved for fixed $\gamma>0$.
\end{enumerate}

\section{Further directions}

The framework suggests several follow-up problems that are best treated separately from the present note.

\paragraph{Restricted circuit models.}
Determine the amortization curve for formulas, bounded fan-out circuits, monotone circuits, $AC^0$, threshold circuits, and branching programs.  The affine-line construction relies on routing, XOR-like linear reconstruction, and recursive fan-out; restricted models may expose clean anti-sharing theorems.

\paragraph{Conditional-choice classifications.}
Classify finite Boolean bases by the efficiency with which they can implement conditional resource selection.  Monotone bases lack ordinary muxes, while functionally complete bases implement them at constant cost.  A Post-lattice classification of direct-sum versus mass-production behavior would make the role of conditional choice precise.

\paragraph{Quantum mass production.}
Kretschmer establishes the classical Uhlig range for arbitrary quantum states and unitaries \cite{Kretschmer2023}.  The present proof uses finite-field linear encoding of Boolean restrictions and does not transfer directly.  It nevertheless raises the concrete question of whether a quantum analogue of grouped congestion and recursive reuse can reach $2^{\gamma n}$ copies.

\paragraph{Sharing complexity as a separate invariant.}
The gap between $\C(f)$ and the amount of nonshareable work in $f^{\times t}$ can be exponential.  It may be fruitful to define a ``sharing rank'' or ``congestion profile'' of a circuit representation, rather than treating circuit size as the only invariant.

\appendix

\section{Parameter inequalities for the bootstrap}

For reference, the bootstrap uses only the following open interval.  Assuming
\[
 \gamma<\frac1{2-\beta},
\]
we have
\[
 2\gamma-1<\beta\gamma.
\]
Choose
\[
 c\in(2\gamma-1,\beta\gamma).
\]
For sufficiently large fixed $\ell$,
\[
 c>2\gamma-1+\frac{2(1-\gamma)}\ell.
\]
The scheduling condition follows.  The three overhead exponents are then bounded as follows:
\begin{align*}
 2\gamma-c+\frac{1-\gamma}{\ell}
 &<1-\frac{1-\gamma}{\ell}<1,\\
 \gamma+\frac{1-\gamma}{\ell}
 &<1,\\
 c+1-\gamma
 &<1,
\end{align*}
where the last line uses $c<\beta\gamma\le\gamma$.  For the recursive call, choose
\[
 \rho\in(c/\gamma,\beta),
\]
so that $C=2^{cn}\le2^{\rho d}$ with $d=\gamma n+O(1)$.

\section{A compact pseudocode view of the scheduler}

For one request group, the deterministic schedule is conceptually:

\begin{enumerate}
\item Initialize the used-point list $U\leftarrow\varnothing$.
\item For requests $j=1,\ldots,g$:
  \begin{enumerate}
  \item For every $y\in U\setminus\{u_j\}$, compute the canonical projective rank of $y-u_j$.
  \item Sort and deduplicate the forbidden ranks.
  \item Select the least rank in $[D_q]$ not present.
  \item Enumerate the $q-1$ non-target points on the corresponding line.
  \item Append those points to $U$.
  \end{enumerate}
\end{enumerate}

The sequential presentation describes a combinational circuit obtained by unrolling all steps.  Its size is the sum of the stage sizes, not its running time on a RAM.

\section{Original versus sparse Hirahara parameters}

\begin{center}
\renewcommand{\arraystretch}{1.25}
\begin{tabular}{>{\raggedright\arraybackslash}p{0.34\textwidth}cc}
\toprule
Quantity & Original layout & Sparse regular layout\\
\midrule
Masked share bits per formula & $m\Delta\le\Delta^2$ & $L(\varphi)\le\Delta$\\
Target input length & $O(\log n_0+\Delta^2)$ & $O(\log n_0+\Delta)$\\
Extraction dependence $m'$ & $m\Delta\le\Delta^2$ & $L(\varphi)\le\Delta$\\
Amplification advantage $\eps$ & $\eps_0/(2\Delta^2)$ & $\eps_0/(2\Delta)$\\
Calls to one $f_k$ & $O(\Delta^3\log n_0/\eps_0)$ & $O(\Delta^2\log n_0/\eps_0)$\\
Support density before padding & at least scale $2^{-\Delta^2}$ & exactly $2^{-\Delta}$\\
Convenient degree choice & $\Theta(\sqrt{\log n_0})$ & $\Theta(\log n_0)$\\
Tracked gap from $g=\Delta^a$ & $(\log N)^{a/2}$ & $(\log N)^a$\\
\bottomrule
\end{tabular}
\end{center}

\begin{thebibliography}{99}

\bibitem{Albrecht1989}
Andreas Albrecht.
\newblock On simultaneous realizations of Boolean functions, with applications.
\newblock In \emph{Parcella '88: Proceedings of the Fourth International Workshop on Parallel Processing by Cellular Automata and Arrays}, pages 51--56. Springer, 1989.
\newblock doi:10.1007/3-540-50647-0\_102.

\bibitem{Beimel2011}
Amos Beimel.
\newblock Secret-sharing schemes: A survey.
\newblock In \emph{Coding and Cryptology}, Lecture Notes in Computer Science 6639, pages 11--46. Springer, 2011.

\bibitem{BL1988}
Josh Cohen Benaloh and Jerry Leichter.
\newblock Generalized secret sharing and monotone functions.
\newblock In \emph{CRYPTO '88}, pages 27--35. Springer, 1988.

\bibitem{Hirahara2022}
Shuichi Hirahara.
\newblock NP-hardness of learning programs and partial MCSP.
\newblock \emph{Electronic Colloquium on Computational Complexity}, Report TR22-119, 2022.

\bibitem{IKOS2004}
Yuval Ishai, Eyal Kushilevitz, Rafail Ostrovsky, and Amit Sahai.
\newblock Batch codes and their applications.
\newblock In \emph{Proceedings of STOC 2004}, pages 262--271. ACM, 2004.

\bibitem{ISN1993}
Mitsuru Ito, Akira Saito, and Takao Nishizeki.
\newblock Multiple assignment scheme for sharing secret.
\newblock \emph{Journal of Cryptology}, 6(1):15--20, 1993.

\bibitem{Kretschmer2023}
William Kretschmer.
\newblock Quantum mass production theorems.
\newblock In \emph{18th Conference on the Theory of Quantum Computation, Communication and Cryptography}, LIPIcs 266, Article 10, 2023.

\bibitem{Lupanov1958}
Oleg B. Lupanov.
\newblock A method of circuit synthesis.
\newblock \emph{Izvestiya VUZ, Radiofizika}, 1:120--140, 1958.  In Russian.

\bibitem{NW1994}
Noam Nisan and Avi Wigderson.
\newblock Hardness vs. randomness.
\newblock \emph{Journal of Computer and System Sciences}, 49(2):149--167, 1994.

\bibitem{Paul1976}
Wolfgang J. Paul.
\newblock Realizing Boolean functions on disjoint sets of variables.
\newblock \emph{Theoretical Computer Science}, 2(3):383--396, 1976.
\newblock doi:10.1016/0304-3975(76)90089-X.

\bibitem{Shannon1949}
Claude E. Shannon.
\newblock The synthesis of two-terminal switching circuits.
\newblock \emph{Bell System Technical Journal}, 28(1):59--98, 1949.

\bibitem{Uhlig1974}
Dietmar Uhlig.
\newblock On the synthesis of self-correcting schemes from functional elements with a small number of reliable elements.
\newblock \emph{Mathematical Notes of the Academy of Sciences of the USSR}, 15(6):558--562, 1974.

\bibitem{Uhlig1992}
Dietmar Uhlig.
\newblock Networks computing Boolean functions for multiple input values.
\newblock In M. S. Paterson, editor, \emph{Boolean Function Complexity}, London Mathematical Society Lecture Note Series, pages 165--173. Cambridge University Press, 1992.

\bibitem{Wegener1987}
Ingo Wegener.
\newblock \emph{The Complexity of Boolean Functions}.
\newblock Wiley-Teubner, 1987.

\end{thebibliography}

\end{document}
